\documentclass[10pt,a4paper]{amsart}
\usepackage[margin=19mm]{geometry}
\usepackage{amsmath,amssymb,mathtools}
\usepackage[scr=rsfs,cal=euler]{mathalfa}
\usepackage{xfrac}
\usepackage[unicode,hidelinks,bookmarks=false]{hyperref}
\newcommand{\ie}{i.e.\ }
\newcommand{\cf}{cf.\ }
\newcommand{\resp}{respectively\ }
\newcommand{\Subsection}{\subsection}
\providecommand{\nobreakdash}{\nobreak\hbox{-}}
\setlength{\emergencystretch}{2em}
\multlinegap0pt
\usepackage{bm}
\usepackage{bbm}
\usepackage{dsfont}
\usepackage{xspace}
\usepackage{cancel}
\usepackage{mathtools}
\usepackage{amsthm}
\makeatletter
\@ifundefined{theorem}{\newtheorem{theorem}{Theorem}}{}
\@ifundefined{lemma}{\newtheorem{lemma}{Lemma}}{}
\makeatother
\title[Optimal nodal control and the exponential turnpike property ]{Optimal nodal control and the exponential turnpike property for the flow in gas networks}
\author{Martin Gugat, Michael Herty, Michael Schuster, Yizhou Zhou}
\date{Source: arXiv:2609.26572v1 (CC BY 4.0)}
\begin{document}
\maketitle

\noindent\emph{Source and licence.} Optimal nodal control and the exponential turnpike property for the flow in gas networks. Version arXiv:2609.26572v1 (CC BY 4.0), distributed under the Creative Commons Attribution 4.0 International licence, as recorded in the arXiv licence metadata for this exact version and shown on the abstract page https://arxiv.org/abs/2609.26572v1. Full rights notice: licenses/arxiv-2609.26572-CC-BY-4.0.txt.\par\smallskip

\noindent\textit{Editorial scope.} Counted unit: the Lemma whose source label is \texttt{lemma3.2-eq1} in the article (source0.tex lines 581--587, source label \texttt{lemma3.2-eq1}), delivered together with the article's own complete original proof of it (source0.tex lines 588--830, one \texttt{proof} environment of 243 lines). No mathematics is written, repaired or re-derived: statement and proof are the article's own text line for line. Required context retained uncounted: eq1.11 (lines 326--331); eq1.12 (lines 326--331); eq1.13 (lines 333--336). Every definition, display and label of those lines is retained unchanged. Omitted from this package and not redistributed: 1--325, 332--332, 337--580, 831--1454. Nothing inside a retained excerpt is omitted or altered: every retained body is delivered exactly as the article's lines are written, so no omission receipt of any kind accompanies this package. Citations: the retained excerpts carry no citation command at all, so no bibliography entry of the article is transcribed and no reference is redistributed. Reference adaptation: none was needed and none was made; every reference inside the delivered unit resolves inside it, so no unresolved reference or citation is produced. Printed slips kept literal: no printed word, symbol, hypothesis or number of the counted statement or of its proof was altered. Layout and class: the article's own class is \texttt{article} with packages including amssymb, amsthm, hyperref, amsmath, bm, amsfonts, arydshln, verbatim, graphicx, float, subfigure, multirow, color, appendix; the wrapper is the lane's pinned \texttt{amsart} editorial wrapper at 10pt on A4, which restates the theorem-like environments the retained text needs and adds the article's own macro definitions that the retained text needs (none), with extra wrapper packages (bm, bbm, dsfont, xspace, cancel, mathtools). The delivered numbering is the unit's own. Assets: no retained span contains a figure environment, a graphics-inclusion command, a PDF-page inclusion, a table or a TikZ picture; no image file is copied into this package, the package declares no asset dependency and no image file is redistributed with it. Licence: version arXiv:2609.26572v1 of this article is distributed under the Creative Commons Attribution 4.0 International licence, as recorded in the arXiv licence metadata for this exact version and shown on the abstract page https://arxiv.org/abs/2609.26572v1. A full rights notice is delivered at licenses/arxiv-2609.26572-CC-BY-4.0.txt. Characters: every retained excerpt is pure ASCII, so the delivered engine, XeLaTeX with its default Latin Modern text font, reports no missing character.
\begin{align}
    r_-^{(1)}(L,t)&=\Phi_1\left({r}_+^{(1)}(L,t)+\bar{R}_+^{(1)}(L),{r}_-^{(2)}(0,t)+\bar{R}_-^{(2)}(0),u_0(t)\right)\notag \\[1mm]
    &\quad -\Phi_1(\bar{R}_+^{(1)}(L),\bar{R}_-^{(2)}(0),\bar{u}_0) \label{eq1.11}\\[2mm]
    r_+^{(2)}(0,t)&=\Phi_2\left({r}_+^{(1)}(L,t)+\bar{R}_+^{(1)}(L),{r}_-^{(2)}(0,t)+\bar{R}_-^{(2)}(0),u_0(t)\right)\notag \\[1mm]
    &\quad -\Phi_2(\bar{R}_+^{(1)}(L),\bar{R}_-^{(2)}(0),\bar{u}_0)  \label{eq1.12}
\end{align}

\begin{align}\label{eq1.13}
    r_+^{(1)}(0,t)=u_1(t)-\bar{u}_1,\qquad
    r_-^{(2)}(L,t)=u_2(t)-\bar{u}_2.
\end{align}

\begin{lemma}
For sufficiently small $L$, there exist constants $A^{(i)}$, $B^{(i)}$, $C^{(i)}$ and $D^{(i)}$ such that the following relation holds 
\begin{align}\label{lemma3.2-eq1} 
\sum_{j=1}^2\left(I^{(j)}_0+I^{(j)}_L\right)\leq C\sum_{i=0}^2\Big(|u_i(t)-\bar{u}_i|^2+|\partial_tu_i(t)|^2+|\partial_t^2u_i(t)|^2\Big).
\end{align}
Here $C>0$ is a generic constant.
\end{lemma}

\begin{proof}
Recall that the boundary conditions of $r^{(i)}(x,t)$ are given by \eqref{eq1.11}--\eqref{eq1.13}. We take time derivatives on these relations to obtain
\begin{align}
    p_-^{(1)}(L,t)&=\partial_{R_+}\Phi_1\cdot p_+^{(1)}(L,t) + \partial_{R_-}\Phi_1\cdot p_-^{(2)}(0,t)  +\partial_{u}\Phi_1\cdot \partial_tu_0(t)\label{eq-p1}\\[2mm]
    p_+^{(2)}(0,t)&=\partial_{R_+}\Phi_2\cdot p_+^{(1)}(L,t) + \partial_{R_-}\Phi_2\cdot p_-^{(2)}(0,t)  +\partial_{u}\Phi_2\cdot \partial_tu_0(t)\label{eq-p2}
\end{align}
and
\begin{align}
    p_+^{(1)}(0,t)=\partial_tu_1(t),\qquad
    p_-^{(2)}(L,t)=\partial_tu_2(t).
\end{align}
Moreover, differentiating the coupling conditions twice with
respect to time, we obtain
\begin{align}
s_-^{(1)}(L,t)
={}&
\partial_{R_+}\Phi_1\,s_+^{(1)}(L,t)
+\partial_{R_-}\Phi_1\,s_-^{(2)}(0,t)
+\partial_u\Phi_1\,\partial_t^2u_0(t)
\nonumber\\
&+
\partial_{R_+R_+}\Phi_1
\left|p_+^{(1)}(L,t)\right|^2
+
\partial_{R_-R_-}\Phi_1
\left|p_-^{(2)}(0,t)\right|^2
+
\partial_{uu}\Phi_1
\left|\partial_tu_0(t)\right|^2
\nonumber\\
&+
2\partial_{R_+R_-}\Phi_1
p_+^{(1)}(L,t)p_-^{(2)}(0,t)
\nonumber\\
&+
2\partial_{R_+u}\Phi_1
p_+^{(1)}(L,t)\partial_tu_0(t)
+
2\partial_{R_-u}\Phi_1
p_-^{(2)}(0,t)\partial_tu_0(t),
\label{eq-s1}
\end{align}
and
\begin{align}
s_+^{(2)}(0,t)
={}&
\partial_{R_+}\Phi_2\,s_+^{(1)}(L,t)
+\partial_{R_-}\Phi_2\,s_-^{(2)}(0,t)
+\partial_u\Phi_2\,\partial_t^2u_0(t)
\nonumber\\
&+
\partial_{R_+R_+}\Phi_2
\left|p_+^{(1)}(L,t)\right|^2
+
\partial_{R_-R_-}\Phi_2
\left|p_-^{(2)}(0,t)\right|^2
+
\partial_{uu}\Phi_2
\left|\partial_tu_0(t)\right|^2
\nonumber\\
&+
2\partial_{R_+R_-}\Phi_2
p_+^{(1)}(L,t)p_-^{(2)}(0,t)
\nonumber\\
&+
2\partial_{R_+u}\Phi_2
p_+^{(1)}(L,t)\partial_tu_0(t)
+
2\partial_{R_-u}\Phi_2
p_-^{(2)}(0,t)\partial_tu_0(t).
\label{eq-s2}
\end{align}
% Moreover, we also get
% \begin{align}
%     s_-^{(1)}(L,t)&=\partial_{R_+}\Phi_1\cdot s_+^{(1)}(L,t) + \partial_{R_-}\Phi_1\cdot s_-^{(2)}(0,t)  +\partial_{u}\Phi_1\cdot \partial_t^2u_0(t)\nonumber\\[2mm]
%     &~+\partial_{R_+}^2\Phi_1\cdot |p_+^{(1)}(L,t)|^2 + \partial_{R_-}^2\Phi_1\cdot |p_-^{(2)}(0,t)|^2  +\partial_{u}^2\Phi_1\cdot |\partial_tu_0(t)|^2 \label{eq-s1}\\[2mm]
%     s_+^{(2)}(0,t)&=\partial_{R_+}\Phi_2\cdot s_+^{(1)}(L,t) + \partial_{R_-}\Phi_2\cdot s_-^{(2)}(0,t)  +\partial_{u}\Phi_2\cdot \partial_t^2u_0(t)\nonumber\\[2mm]
%     &~+\partial_{R_+}^2\Phi_2\cdot |p_+^{(1)}(L,t)|^2 + \partial_{R_-}^2\Phi_2\cdot |p_-^{(2)}(0,t)|^2  +\partial_{u}^2\Phi_2\cdot |\partial_tu_0(t)|^2\label{eq-s2}
% \end{align}
and
\begin{align}
    s_+^{(1)}(0,t)=\partial_t^2u_1(t),\qquad
    s_-^{(2)}(L,t)=\partial_t^2u_2(t).
\end{align}
Since $\Phi_1$ and $\Phi_2$ are of class $C^2$ and their
arguments remain in a compact set, all first- and second-order
derivatives of $\Phi_1$ and $\Phi_2$ are uniformly bounded.
Moreover, the smallness assumptions on the solution and controls,
together with the equations, imply that
\[
|p_+^{(1)}(L,t)|
+|p_-^{(2)}(0,t)|
+|\partial_tu_0(t)|
\leq \delta
\]
for some small constant $\delta>0$. Thus,
\[
|p_+^{(1)}(L,t)|^4
\leq
\delta^2|p_+^{(1)}(L,t)|^2,
\]
and analogous estimates hold for
$p_-^{(2)}(0,t)$ and $\partial_tu_0(t)$. The mixed terms are
controlled by $2|ab|\leq |a|^2+|b|^2$. Hence, there exists a constant $\widehat C_0>0$  such that
\begin{align}
|s_-^{(1)}(L,t)|^2+|s_+^{(2)}(0,t)|^2
\leq{}&
\widehat C_0\Big(
|s_+^{(1)}(L,t)|^2
+|s_-^{(2)}(0,t)|^2
+|\partial_t^2u_0(t)|^2
\nonumber\\
&\quad
+|p_+^{(1)}(L,t)|^2
+|p_-^{(2)}(0,t)|^2
+|\partial_tu_0(t)|^2
\Big).
\label{I_s-eq2}
\end{align}
We add all terms in $\sum_{j=1}^2(I_0^{(j)}+I_L^{(j)})$ and reorganize them as
$$
\sum_{j=1}^2(I_0^{(j)}+I_L^{(j)}) = I_{r}+I_p+I_s,
$$
where $I_r$ contains the terms $r_{\pm}^{(i)}$:
$$
I_r=\sum_{i=1}^2\Big[A^{(i)} |r_+^{(i)}(0,t)|^2-B^{(i)} |r_-^{(i)}(0,t)|^2
-A^{(i)}h_+^{(i)}(L)|r_+^{(i)}(L,t)|^2+B^{(i)}h_-^{(i)}(L)|r_-^{(i)}(L,t)|^2\Big],
$$
while $I_p$ and $I_s$ contain the terms $p_{\pm}^{(i)}$ and $s_{\pm}^{(i)}$ respectively:
\begin{align*}
I_p = &~\sum_{i=1}^2C^{(i)} \widehat{\lambda}^{(i)}_+(0) |p_+^{(i)}(0,t)|^2 -D^{(i)}\widehat{\lambda}^{(i)}_-(0) |p_-^{(i)}(0,t)|^2 -C^{(i)}h_+^{(i)}(L)\widehat{\lambda}^{(i)}_+(L) |p_+^{(i)}(L,t)|^2 \\[2mm]
&~+ D^{(i)}h_-^{(i)}(L)\widehat{\lambda}^{(i)}_-(L)|p_-^{(i)}(L,t)|^2\\[2mm]
I_s = &~\sum_{i=1}^2C^{(i)} \widehat{\lambda}^{(i)}_+(0) |s_+^{(i)}(0,t)|^2 -D^{(i)}\widehat{\lambda}^{(i)}_-(0) |s_-^{(i)}(0,t)|^2 -C^{(i)}h_+^{(i)}(L)\widehat{\lambda}^{(i)}_+(L) |s_+^{(i)}(L,t)|^2 \\[2mm]
&~+ D^{(i)}h_-^{(i)}(L)\widehat{\lambda}^{(i)}_-(L)|s_-^{(i)}(L,t)|^2. 
\end{align*}
Here the coefficients are defined by
\begin{align*}
&\widehat{\lambda}^{(i)}_+(0)  =\bar{\lambda}^{(i)}_+(0)-\frac{r^{(i)}_+(0,t)+r^{(i)}_-(0,t)}{2},\qquad 
\widehat{\lambda}^{(i)}_-(0)=|\bar{\lambda}^{(i)}_-(0)|+\frac{r^{(i)}_+(0,t)+r^{(i)}_-(0,t)}{2}\\[2mm]
&\widehat{\lambda}^{(i)}_+(L)  =\bar{\lambda}^{(i)}_+(L)+\frac{r^{(i)}_+(L,t)+r^{(i)}_-(L,t)}{2},\qquad 
\widehat{\lambda}^{(i)}_-(L)=|\bar{\lambda}^{(i)}_-(L)|-\frac{r^{(i)}_+(L,t)+r^{(i)}_-(L,t)}{2}.
\end{align*}
For sufficiently small $\epsilon(T)$, we know that these four numbers are all positive.

\textbf{Estimate for $I_r$:} For $r_\pm^{(i)}$ terms, we use \eqref{eq1.13} and compute  
\begin{align}
I_r
% =&\sum_{i=1}^2\Big[A^{(i)} |r_+^{(i)}(0,t)|^2-B^{(i)} |r_-^{(i)}(0,t)|^2-A^{(i)}h_+^{(i)}(L)|r_+^{(i)}(L,t)|^2+B^{(i)}h_-^{(i)}(L)|r_-^{(i)}(L,t)|^2\Big]\\[2mm]
\leq&~ A^{(1)} |u_1(t)-\bar{u}_1|^2 
-A^{(1)}h_+^{(1)}(L)|r_+^{(1)}(L,t)|^2+B^{(1)}h_-^{(1)}(L)|r_-^{(1)}(L,t)|^2 \nonumber\\[2mm]
&+A^{(2)} |r_+^{(2)}(0,t)|^2-B^{(2)} |r_-^{(2)}(0,t)|^2
 +B^{(2)}h_-^{(2)}(L)|u_2(t)-\bar{u}_2|^2\nonumber \\[2mm]
 \leq& -C_1|\alpha|^2+C_2|\beta|^2 +C_3\big(|u_1(t)-\bar{u}_1|^2+|u_2(t)-\bar{u}_2|^2\big)\label{I_r-eq1}
\end{align}
Here we define $$\alpha=(r_+^{(1)}(L,t), r_-^{(2)}(0,t)),\qquad \beta=(r_-^{(1)}(L,t), r_+^{(2)}(0,t)).$$
The constants are given by  
\begin{equation}\label{I_r-eq2}
C_1=\min\{A^{(1)}h_+^{(1)}(L), B^{(2)} \} ,\quad C_2=\max\{B^{(1)}h_-^{(1)}(L), A^{(2)} \}
\end{equation}
and $C_3=\max\{A^{(1)}, B^{(2)}h_-^{(2)}(L)\}$.
From \eqref{eq1.11} and \eqref{eq1.12}, we can use the local Lipschitz continuity to conclude that
% $$
% \beta=\Psi(\alpha,u_0(t))
% $$
% To solve this nonlinear equation and get $\beta=\beta(\alpha,u_0)$ near the point $(\alpha,\beta,u_0)=(0,0,\bar{u}_0)$, we compute
% $$
% \frac{\partial \beta}{\partial \alpha} = \left[\frac{\partial \Psi}{\partial \beta}\right]^{-1}\frac{\partial \Psi}{\partial \alpha},\qquad 
% \frac{\partial \beta}{\partial u_0} = \left[\frac{\partial \Psi}{\partial \beta}\right]^{-1}\frac{\partial \Psi}{\partial u_0}.
% $$
% Since $\Psi$ is $C^2$-smooth and we restrict our attention to solutions lying in a compact set, it follows that
% $$
% \left|\frac{\partial \beta}{\partial \alpha}\right|\leq C,\qquad \left|\frac{\partial \beta}{\partial u_0}\right|\leq C.
% $$
% Thus there exists a constant $C_0>0$ such that
\begin{equation}\label{beta-au}
|\beta|^2\leq C_0(|\alpha|^2+|u_0-\bar{u}_0|^2)
\end{equation}
for $(\alpha,\beta,u_0)$ sufficiently close to $(0,0,\bar{u}_0)$ with $C_0>0$ a constant.
Substituting this into \eqref{I_r-eq1}, we have 
\begin{align*}
    I_r &\leq (C_2C_0-C_1)|\alpha|^2+C_2C_0 |u_0-\bar{u}_0|^2+C_3\big(|u_1(t)-\bar{u}_1|^2+|u_2(t)-\bar{u}_2|^2\big).
\end{align*}
According to the expression in \eqref{I_r-eq2}, we take $A^{(1)}$, $B^{(2)}$ sufficiently large and $A^{(2)}$, $B^{(1)}$ sufficiently small such that $C_2C_0-C_1$ is negative. Consequently, we obtain the estimate for $r_\pm^{(i)}$ terms.

\textbf{Estimate for $I_p$ and $I_s$:} For $p_\pm^{(i)}$ terms, we compute
% \begin{align*}
% I^{(p)} = &~\sum_{i=1}^2C^{(i)} \widehat{\lambda}^{(i)}_+(0) |p_+^{(i)}(0,t)|^2 -D^{(i)}\widehat{\lambda}^{(i)}_-(0) |p_-^{(i)}(0,t)|^2 -C^{(i)}h_+^{(i)}(L)\widehat{\lambda}^{(i)}_+(L) |p_+^{(i)}(L,t)|^2 \\[2mm]
% &~+ D^{(i)}h_-^{(i)}(L)\widehat{\lambda}^{(i)}_-(L)|p_-^{(i)}(L,t)|^2.
% \end{align*}
% Then it follows that 
\begin{align}
I_p \leq &~ C^{(1)} \widehat{\lambda}^{(1)}_+(0) |\partial_t u_1(t)|^2  -C^{(1)}h_+^{(1)}(L)\widehat{\lambda}^{(1)}_+(L) |p_+^{(1)}(L,t)|^2 + D^{(1)}h_-^{(1)}(L)\widehat{\lambda}^{(1)}_-(L)|p_-^{(1)}(L,t)|^2 \nonumber \\[3mm]
&+ C^{(2)} \widehat{\lambda}^{(2)}_+(0) |p_+^{(2)}(0,t)|^2 -D^{(2)}\widehat{\lambda}^{(2)}_-(0) |p_-^{(2)}(0,t)|^2 + D^{(2)}h_-^{(2)}(L)\widehat{\lambda}^{(2)}_-(L)|\partial_tu_2(t)|^2 \nonumber \\[2mm]
\leq & -\bar{C}_1\Big(|p_+^{(1)}(L,t)|^2 + |p_-^{(2)}(0,t)|^2 \Big)+\bar{C}_2\Big(|p_-^{(1)}(L,t)|^2+
|p_+^{(2)}(0,t)|^2\Big)\nonumber  \\[2mm]
&+\bar{C}_3\Big(|\partial_tu_1(t)|^2+|\partial_tu_2(t)|^2\Big)\label{I_p-eq1}
\end{align}
with
\begin{align*}
\bar{C}_1&=\min\{C^{(1)}h_+^{(1)}(L)\widehat{\lambda}^{(1)}_+(L) ,~D^{(2)}\widehat{\lambda}^{(2)}_-(0)\}>0,\\[2mm] 
\bar{C}_2&=\max\{D^{(1)}h_-^{(1)}(L)\widehat{\lambda}^{(1)}_-(L) ,~C^{(2)}\widehat{\lambda}^{(2)}_+(0)\}>0,\\[2mm]
\bar{C}_3&=\max\{C^{(1)} \widehat{\lambda}^{(1)}_+(0) ,~D^{(2)}h_-^{(2)}(L)\widehat{\lambda}^{(2)}_-(L)\}>0.
\end{align*}
In \eqref{eq-p1}--\eqref{eq-p2}, $\Phi_1$ and $\Phi_2$ are $C^2$-smooth functions with respect to the components and the solution remains in a compact set. Thus we know that there exists a constant $\bar{C}_0>0$ such that
% $$
% \max\{|\partial_{R_+}\Phi_k|,|\partial_{R_-}\Phi_k|, |\partial_{u}\Phi_k|,|\partial_{R_+}^2\Phi_k|,|\partial_{R_-}^2\Phi_k|, |\partial_{u}^2\Phi_k|\}\leq \bar{C}_0
% $$
% for $k=1,2$. Therefore, we derive from \eqref{eq-p1}--\eqref{eq-p2} that 
\begin{align}\label{I_p-eq2}
    |p_-^{(1)}(L,t)|^2+|p_+^{(2)}(0,t)|^2&\leq \bar{C}_0 \Big(|p_+^{(1)}(L,t)|^2 + |p_-^{(2)}(0,t)|^2  + |\partial_tu_0(t)|^2\Big).
\end{align}

Similarly, for $s_\pm^{(i)}$ terms, we have 
\begin{align}
I_s
% = &~\sum_{i=1}^2C^{(i)} \widehat{\lambda}^{(i)}_+(0) |s_+^{(i)}(0,t)|^2 -D^{(i)}\widehat{\lambda}^{(i)}_-(0) |s_-^{(i)}(0,t)|^2 -C^{(i)}h_+^{(i)}(L)\widehat{\lambda}^{(i)}_+(L) |s_+^{(i)}(L,t)|^2 \\[2mm]
% &~+ D^{(i)}h_-^{(i)}(L)\widehat{\lambda}^{(i)}_-(L)|s_-^{(i)}(L,t)|^2 \\[2mm]
\leq &~ C^{(1)} \widehat{\lambda}^{(1)}_+(0) |\partial_t^2 u_1(t)|^2  -C^{(1)}h_+^{(1)}(L)\widehat{\lambda}^{(1)}_+(L) |s_+^{(1)}(L,t)|^2 + D^{(1)}h_-^{(1)}(L)\widehat{\lambda}^{(1)}_-(L)|s_-^{(1)}(L,t)|^2\nonumber\\[3mm]
&+ C^{(2)} \widehat{\lambda}^{(2)}_+(0) |s_+^{(2)}(0,t)|^2 -D^{(2)}\widehat{\lambda}^{(2)}_-(0) |s_-^{(2)}(0,t)|^2 + D^{(2)}h_-^{(2)}(L)\widehat{\lambda}^{(2)}_-(L)|\partial_t^2u_2(t)|^2\nonumber\\[2mm]
\leq & -\bar{C}_1\Big(|s_+^{(1)}(L,t)|^2 + |s_-^{(2)}(0,t)|^2 \Big)+\bar{C}_2\Big(|s_-^{(1)}(L,t)|^2+
|s_+^{(2)}(0,t)|^2\Big)\nonumber\\[2mm]
&+\bar{C}_3\Big(|\partial_t^2u_1(t)|^2+|\partial_t^2u_2(t)|^2\Big).\label{I_s-eq1}
\end{align}
% By \eqref{eq-s1}--\eqref{eq-s2}, we know that there exists a constant $\widehat{C}_0>0$ such that
% \begin{align}
%     |s_-^{(1)}(L,t)|^2+|s_+^{(2)}(0,t)|^2\leq &~\widehat{C}_0\Big(|s_+^{(1)}(L,t)|^2 + |s_-^{(2)}(0,t)|^2  + |\partial_t^2u_0(t)|^2\nonumber \\[2mm]
%     &+|p_+^{(1)}(L,t)|^2 + |p_-^{(2)}(0,t)|^2  + |\partial_tu_0(t)|^2\Big).\label{I_s-eq2}
% \end{align}
Combining \eqref{I_s-eq2},\eqref{I_p-eq1}--\eqref{I_s-eq1},  we have 
\begin{align*}
I_p +I_s
\leq & ~(\widehat{C}_0\bar{C}_2-\bar{C}_1)\Big(|s_+^{(1)}(L,t)|^2 + |s_-^{(2)}(0,t)|^2 \Big) \\[2mm]
&+(\widehat{C}_0\bar{C}_2+\bar{C}_0\bar{C}_2-\bar{C}_1)\Big(|p_+^{(1)}(L,t)|^2 + |p_-^{(2)}(0,t)|^2 \Big)\\[2mm]
&+\bar{C}_3\Big(|\partial_tu_1(t)|^2+|\partial_tu_2(t)|^2+|\partial_t^2u_1(t)|^2+|\partial_t^2u_2(t)|^2\Big)\\[2mm]
&+\bar{C}_2(\bar{C}_0+\widehat{C}_0)\Big(|\partial_tu_0(t)|^2+|\partial_t^2u_0(t)|^2\Big).
\end{align*}
According to the expression of $\bar{C}_1$ and $\bar{C}_2$, we can choose $C^{(1)}$, $D^{(2)}$ to be sufficiently large and $C^{(2)}$, $D^{(1)}$ to be sufficiently small such that 
$$
\widehat{C}_0\bar{C}_2-\bar{C}_1 <
\widehat{C}_0\bar{C}_2+\bar{C}_0\bar{C}_2-\bar{C}_1 <0.
$$
This gives the estimate for $I_p$ and $I_s$. We complete the proof.
\end{proof}

\end{document}
